\documentclass[10pt,twocolumn]{article}
\usepackage[T1]{fontenc}
\usepackage[utf8]{inputenc}
\usepackage{lmodern}
\usepackage[margin=1.8cm,top=2.4cm,bottom=2cm,columnsep=0.6cm]{geometry}
\usepackage{fancyhdr}
\usepackage{titlesec}
\usepackage{booktabs}
\usepackage{amsmath,amssymb,amsfonts}
\usepackage{microtype}
\usepackage{xcolor}
\usepackage{mdframed}
\usepackage{parskip}
\usepackage{enumitem}
\usepackage{hyperref}

\definecolor{rulered}{RGB}{180,30,30}
\definecolor{boxgray}{RGB}{245,245,245}
\definecolor{keywordgray}{RGB}{100,100,100}

\pagestyle{fancy}
\fancyhf{}
\fancyhead[L]{\textsc{\small Claude Mag}}
\fancyhead[C]{\small\textit{Machine Learning Research Digest}}
\fancyhead[R]{\small 2026-08-31}
\fancyfoot[C]{\small\thepage}
\renewcommand{\headrulewidth}{0.8pt}

\titleformat{\section}{\normalsize\bfseries\color{rulered}}{}{0pt}{}%
  [\vspace{-4pt}\color{rulered}\rule{\columnwidth}{0.4pt}]
\titlespacing{\section}{0pt}{8pt}{4pt}

\hypersetup{colorlinks=true,urlcolor=rulered,linkcolor=black}

\begin{document}
\setlength{\parindent}{0pt}
\setlength{\parskip}{4pt}

{\small\color{keywordgray}\textbf{Keywords:}
  agent memory \textbullet{}
  multimodal retrieval \textbullet{}
  graph neural networks \textbullet{}
  long-horizon agents}\par\vspace{4pt}

{\LARGE\bfseries\raggedright
  GraphMemix: Query-Aware Evidence Forests
  for Long-Term Multimodal Agent Memory}\par\vspace{4pt}

{\large\itshape\raggedright
  A Learned Forest-Construction Policy Adapts Multi-Hop
  Retrieval to Query Intent Without Ground-Truth Path Supervision}\par\vspace{6pt}

{\small\textbf{By} Geng Li, Yuhao Wang, Dong Li,
  Jianye Hao, Yuxin Peng}\par\vspace{2pt}

{\small\color{keywordgray}arXiv:2608.26983 \textbullet{} Preprint, August 2026}
\par\vspace{2pt}\noindent{\color{rulered}\rule{\columnwidth}{0.6pt}}\vspace{6pt}

Long-horizon multimodal agents accumulate heterogeneous memories —
images, video clips, text segments — and must retrieve relevant
evidence across modalities for each new query. Flat dense retrieval
handles heterogeneity but misses multi-hop dependencies; graph memory
captures dependencies but cannot adapt its retrieval structure to
query intent. GraphMemix resolves this by learning a lightweight
policy that dynamically constructs a query-aware evidence forest
from a persistent heterogeneous memory graph, improving retrieval
accuracy by up to $+11.3\%$ over dense retrieval on long-horizon
benchmarks at sub-second query latency.

\begin{mdframed}[backgroundcolor=boxgray,linecolor=rulered,linewidth=1pt,
    innerleftmargin=6pt,innerrightmargin=6pt,
    innertopmargin=6pt,innerbottommargin=6pt]
\textbf{\small AT A GLANCE}\par\vspace{4pt}
\begin{description}[leftmargin=0pt,labelsep=4pt,itemsep=2pt,parsep=0pt,
    font=\small\bfseries]
  \item[Problem:] Multimodal agents need memory retrieval that handles
    cross-modal heterogeneity and multi-hop reasoning simultaneously;
    neither dense retrieval nor static graph schemas do both.
  \item[Why it matters:] As agents operate over longer interaction
    histories with diverse modalities, memory retrieval quality becomes
    a primary bottleneck for task completion.
  \item[Method:] Heterogeneous memory graph with a GNN-based forest
    construction policy that selects which edges to expand, conditioned
    on query embedding; trained with sparse REINFORCE rewards.
  \item[Result:] $+$7.2\% on EgoSchema-Long, $+$11.3\% on MMAgent-Bench,
    $+$6.8 F1 on MultiModal-NarrativeQA vs.\ dense retrieval.
  \item[Evidence:] Removing query-conditioned forest construction (using
    static BFS) degrades by $-$4.3\% on average; cross-modal aggregation
    removes another $-$3.1\%.
\end{description}
\end{mdframed}

\section{The Big Picture}

The memory requirements of AI agents have evolved dramatically.
Early agents had short context windows and static knowledge.
Contemporary multimodal agents are deployed in long-horizon settings
where they accumulate interaction histories spanning hours: images
from visited locations, video clips of observed events, text segments
from conversations and documents.

Effective agents must retrieve relevant evidence from this growing
memory on demand, for queries that are themselves multimodal and
diverse in intent. A query like ``where did I see the red toolbox?''
requires matching a text query to visual memories. A query like
``what sequence of events led to the error in step 7?'' requires
chaining multiple memory entries through causal links.

GraphMemix addresses this joint challenge with a single memory
architecture that adapts its retrieval structure to each query.

\section{Why Existing Approaches Fall Short}

\textbf{Dense retrieval:} Embedding-based similarity search (DPR,
FAISS) retrieves independent top-$k$ nodes from the memory pool.
This handles modality diversity (using shared embedding spaces) but
treats all retrievals as independent — missing multi-hop chains where
the relevant evidence at step 2 depends on what was retrieved at
step 1.

\textbf{Static graph memory:} Systems that maintain a persistent
knowledge graph (MemGraph, SAGE, GAM) capture dependencies between
memory items as typed edges and support multi-hop traversal. However,
their traversal strategy — usually breadth-first or heuristic
depth-first — is fixed regardless of query intent. A depth-first
strategy appropriate for causal chain tracing is suboptimal for
gathering diverse evidence across parallel memory threads.

\textbf{Hybrid systems:} Retrieving from a dense index and then
expanding via graph edges (retrieve-then-expand) partially addresses
this, but the expansion policy is not conditioned on the query,
limiting adaptation.

\section{Core Mechanism}

GraphMemix maintains a \textbf{heterogeneous memory graph}
$G = (V, E)$:

\begin{itemize}[itemsep=2pt,parsep=0pt]
  \item \textbf{Nodes} $V$: memory items typed by modality (text
    segment, image, video clip summary)
  \item \textbf{Edges} $E$: typed relations — temporal (A before B),
    co-occurrence (A and B observed together), causal (A caused B),
    spatial (A located near B), and co-reference (A and B refer
    to the same entity)
\end{itemize}

For each incoming query $q$, the system constructs a
\textbf{query-aware evidence forest} in three steps:

\textbf{Step 1 — Seed retrieval:} A dual-encoder retrieves top-$k$
seed nodes $S \subset V$ by cross-modal similarity between the
query embedding $h_q$ and node embeddings $\{h_v\}$.

\textbf{Step 2 — Forest construction:} A GNN-based policy
$\pi_\phi$ selects which outgoing edges to expand from each seed,
constructing a tree of depth $d$ rooted at each seed. The policy
decision at node $v$ for edge $(v, u)$ is:
\[
  P_\phi(u | v, q) = \sigma\left(W_\phi [h_v \| h_u \| h_q \| r_{vu}]\right)
\]
where $r_{vu}$ is the relation type embedding. The top-$b$ expansions
(branch factor) are selected; the policy is applied recursively for
$d$ layers.

\textbf{Step 3 — Evidence aggregation:} Nodes across all trees
(deduplicated) are encoded by a cross-modal transformer. A relevance
scorer ranks them, and the top-$m$ are passed to the agent.

\section{Technical Formulation}

The forest for query $q$ from seed set $S$ is:
\[
  \mathcal{F}(q, S) = \bigcup_{s \in S} T_\phi(s, q, d, b)
\]
where $T_\phi$ is the tree built by the policy from seed $s$ with
depth $d$ and branch factor $b$. Complexity is
$O(|S| \cdot b^d \cdot c_\phi)$ where $c_\phi$ is the GNN forward
pass cost — tractable for $|S|=5, b=3, d=2$ ($\leq$45 policy calls
per query).

The policy is trained with REINFORCE using downstream task accuracy
$R$ as the reward signal and a baseline $\hat{R}$ from dense
retrieval:
\[
  \nabla_\phi \mathcal{J} = \mathbb{E}\left[(R - \hat{R})
    \sum_{e \in \mathcal{F}} \nabla_\phi \log P_\phi(e)\right]
\]

No ground-truth paths through the graph are required; the policy
learns which edge types to expand under which query conditions from
task-level reward alone.

The cross-modal aggregation transformer uses cross-attention between
text and non-text nodes at each depth level, with modality-type
embeddings injected via learned positional offsets.

\section{Learning or Inference Procedure}

\textbf{Training:}
\begin{enumerate}[itemsep=2pt,parsep=0pt]
  \item Pre-train seed encoder with contrastive cross-modal loss
  \item Train forest policy with REINFORCE on each benchmark,
    using 80\% of instances for training
  \item Fine-tune cross-modal aggregation transformer on evidence
    re-ranking objective
\end{enumerate}

\textbf{Inference:} Given query $q$, retrieve seeds ($\sim$10ms),
run forest construction policy ($\sim$50ms for $|S|=5, b=3, d=2$
on an A100), aggregate and rank evidence ($\sim$270ms). Total:
$\sim$0.32s per query for a 10K-node graph.

\section{What the Guarantee Says}

No formal convergence guarantees are provided for the REINFORCE
training. Empirically, the policy training converges in 3--5
epochs on all three benchmarks. The system provides no worst-case
latency guarantee, but median latency is 0.32s and 99th percentile
is 1.1s on a 10K-node graph.

\section{Experimental Findings}

\begin{itemize}[itemsep=2pt,parsep=0pt]
  \item \textbf{EgoSchema-Long:} $+$7.2\% accuracy vs.\ dense retrieval;
    $+$4.1\% vs.\ best static graph baseline
  \item \textbf{MMAgent-Bench:} $+$11.3\% task completion vs.\ dense
    retrieval; $+$6.8\% vs.\ static graph
  \item \textbf{MultiModal-NarrativeQA:} $+$6.8 F1 vs.\ dense retrieval;
    $+$3.9 F1 vs.\ static graph
  \item \textbf{Latency:} 0.32s median; scales sub-linearly with graph
    size to 100K nodes
\end{itemize}

\section{Ablations and Interpretation}

\textbf{Query-conditioned vs.\ static forest construction:}
Replacing the policy with static BFS expansion degrades by $-$4.3\%
on average — confirming that query-aware adaptation of tree structure
is the primary driver of improvement.

\textbf{Cross-modal aggregation:} Replacing the cross-modal transformer
with modality-specific encoders (no cross-attention) degrades by
$-$3.1\%, confirming that cross-modal evidence fusion adds value
beyond independent retrieval.

\textbf{Depth $d$ and branch factor $b$:} $d{=}2, b{=}3$ is optimal;
$d{=}3$ slightly improves multi-hop tasks but increases latency by
3$\times$. $b{=}1$ (DFS) is best for causal-chain tasks; $b{=}5$
(broad BFS) is best for evidence-gathering tasks; the policy learns
to approximate this per-query adaptively.

\textbf{Reward baseline:} Using dense retrieval performance as
REINFORCE baseline substantially stabilises policy training
compared to a constant baseline (variance reduced by 4.1$\times$).

\section*{Reference}

Geng Li, Yuhao Wang, Dong Li, Jianye Hao, Yuxin Peng.
\textbf{GraphMemix: Query-Aware Evidence Forests for Long-Term
Multimodal Agent Memory.}
arXiv:2608.26983, August 2026.
\url{https://arxiv.org/abs/2608.26983}

\end{document}
